\begin{textsample}{POS Dim 2 – ac – Score 22.00 – ac/ac\_em\_2.txt}  \label{ex:f2_pos_170}
1.
Contexto histórico O \textbf{Currículo} de Referência Único do Acre para o ensino médio foi construído à luz da Base Nacional Comum Curricular - BNCC, documento de caráter normativo que define o conjunto orgânico e progressivo de aprendizagens essenciais que todos os alunos devem desenvolver ao longo das etapas e modalidades da educação básica.
A necessidade de atualização curricular nas unidades federativas pode ser justificada pelos índices comumente identificados nos processos avaliativos de larga escala e indicadores da Educação Básica, como SAEB e IDEB, que mostram desempenhos aquém das metas projetadas para o ensino médio, tanto em nível \textbf{estadual} como em nível nacional desde 2009.
Além disso, a proposta do Novo Ensino Médio também se faz necessária para diminuir a alta evasão dos estudantes, a fim de construir escolas que façam mais sentido para os jovens e que diminuam as desigualdades educacionais.
Após a homologação da BNCC para a etapa do ensino médio, ocorrida em dezembro de 2018, coube às unidades federativas a articulação de uma equipe que iria, então, elaborar um novo \textbf{currículo}, seguindo todas as orientações que foram estipuladas pelo MEC e seus parceiros.
O Conselho Nacional de \textbf{Secretários} de Educação – CONSED, por meio da Frente \textbf{Currículo} e Novo Ensino Médio, realizou diversos encontros com as equipes responsáveis pela escrita dos \textbf{currículos} \textbf{estaduais}, elaborando e disponibilizando materiais norteadores para a atividade de escrita.
Dessa forma, os \textbf{currículos} são constituídos por uma Formação Geral Básica comum, que, por sua vez, é complementada por uma parte diversificada, que trata da contextualização do ensino, levando em conta as características da região onde o \textbf{currículo} será aplicado.
No estado do Acre, foram realizadas, no início de 2019, pesquisas on-line com os estudantes da \textbf{rede} \textbf{estadual}, envolvendo o nono ano do Ensino Fundamental e a primeira série do ensino médio, para estruturar um modelo que \textbf{atendesse} aos seus anseios.
Desse modo, orientado pelos resultados observados, o \textbf{Currículo} de Referência Único do Acre foi elaborado de modo a retratar, em sua disposição, um documento que considera a realidade do jovem acreano e que se ajusta aos recursos e condições disponíveis em nossa região.
A consulta aos estudantes, pela plataforma Porvir, explicitou direcionamentos para delinear os percursos \textbf{formativos} a serem \textbf{ofertados}, de modo a congregar o equilíbrio entre a autonomia das escolas e, ao mesmo tempo, promover a \textbf{qualidade} e \textbf{equidade} das \textbf{trajetórias} oferecidas, no conjunto das unidades de ensino.
É importante destacar que, no caso do 5º eixo, do Itinerário de Formação \textbf{Técnica} e Profissional, o potencial socioeconômico e ambiental, bem como a capacidade da \textbf{rede} são fundamentos indispensáveis na definição da \textbf{oferta}.
Nesse contexto, a equipe de \textbf{Currículo}, conforme \textbf{Portaria} MEC n{\EmojiFont °} 756 de 03 abril de 2019 que \textbf{institui} o Programa de Apoio à Implementação da Base Nacional Comum Curricular - ProBNCC, elaborou a primeira versão do \textbf{Currículo} de Referência Único do Acre de acordo com a \textbf{reforma} do ensino médio à luz da BNCC, a qual articula possibilidades \textbf{formativas} vinculadas ao projeto de vida dos estudantes.
Durante o processo de escrita também atuaram como colaboradores, outro grupo de docentes da \textbf{rede} pública e privada, \textbf{profissionais} vinculados ao Instituto Federal do Acre – IFAC e à Universidade Federal do Acre – UFAC, constituindo -se uma significativa heterogeneidade de sujeitos ( numa representatividade de sujeitos significativos e heterogêneos ), que contribuíram com proposições de novas metodologias, análise sobre os objetos de conhecimentos abordados, favorecendo o amadurecimento do conceito de ensino e aprendizagem baseado em competências.
Concomitantemente, as escolas-piloto forneceram experiências que indicaram a necessidade de ajustes na matriz curricular, referentes a distribuição da \textbf{carga} \textbf{horária} e disposição dos componentes curriculares na Formação Geral Básica e Itinerário \textbf{Formativo}, \textbf{atendendo} aos modelos de ensino médio em tempo parcial e integral.
Adaptações dos instrumentos de \textbf{rede}, como por exemplo a inserção de ações estratégicas no Plano de Gestão Pedagógico, ajustes da Sequência Didática para \textbf{atender} o trabalho por competências e habilidades.
E novas formas de organização dos tempos e espaços escolares, da gestão do \textbf{currículo}, das abordagens metodológicas, da formação inicial e continuada dos professores, que deverão ser incorporadas nas suas Propostas Pedagógicas e/ou Regimentos Escolares.
Uma vez firmadas as discussões e consolidadas as contribuições agregadas ao texto do \textbf{currículo}, a parte da Formação Geral Básica ( FGB ) entrou em processo de consulta pública, que ocorreu em 2019.
O referido processo se deu, inicialmente, de forma mista, tendo sido realizado um encontro presencial com os professores das áreas de conhecimento, seguido de um período disponibilizado para consulta on-line.
Após esse decurso, o documento da Formação Geral Básica foi \textbf{atualizado}, com base nas contribuições dos professores que participaram e passando também por um processo de seriação das habilidades.
Concluídos os ajustes, nova consulta pública foi realizada, em 2020, apenas via on-line, devido à pandemia do Covid-19, para que fosse submetida à validação dos professores a seriação \textbf{efetuada}.
Em posse das contribuições da segunda consulta, o documento da Formação Geral Básica foi novamente \textbf{atualizado} para sua finalização.
Em paralelo à Consulta Pública da Formação Geral Básica, a equipe de \textbf{Currículo} iniciou a escrita das Rotas de Aprofundamento, parte integrante dos \textbf{Itinerários} \textbf{Formativos}.
Foram elaboradas duas ( 2 ) rotas para cada área de conhecimento, que também foram submetidas ao processo de consulta pública on-line e ajustes mediante as contribuições, concluído o processo, o presente documento foi encaminhado ao Conselho Estadual de Educação do Acre – CEE.
Além das 08 rotas elaboradas e pensando em \textbf{atender} aquele \textbf{município} que tem apenas uma escola de ensino médio, foi elaborada também uma 9º rota no formato integrada de Iniciação Científica e Tecnológica envolvendo unidades das 04 áreas do conhecimento.

% matched lemmas: atender, atualizar, carga, currículo, efetuar, equidade, estadual, formativo, horário, instituir, itinerário, município, oferta, ofertar, portaria, profissional, qualidade, rede, reforma, secretário, trajetória, técnico
\end{textsample}
